\subsection{Complete optimization and recognition measurements}

The final measurements run serially, without tests, audits, builds or
commits overlapping them. Each case has one warm-up round followed by
five measured rounds, each randomly ordering old, old A/A, new and new
A/A endpoints. All 260 measured endpoints and 52 warm-ups complete.
Tables give median primary-arm times and medians of paired ratios;
the latter need not equal ratios of the median times. Each implementation's
certificate hashes remain stable across repetitions, and all comparable
coordinate hashes agree between implementations. Valid matching choices
can make complete certificate hashes differ.

The first experiment includes input construction, span optimization,
an additional independent certificate check and hashing. For the two
source-derived cases, every timed endpoint also constructs the canonical
exterior and solves its primitive cohomology. The binary case scales its
four-crossing cocycle heights by $2^{4096}$ before making the arithmetic
span query; it is not being used as a primitive source-disc certificate.
The random control aggregates 24 complete, independently verified
24-row problems per endpoint to reduce timer noise on small queries.
The other optimizer endpoints contain one query each. This mode completes
140 measured endpoints and 28 warm-ups in 46.753 seconds.

\begin{center}
\small\begin{tabular}{lrrrrr}
\toprule
Input & Old (ms) & New (ms) & Old/new & Old A/A & New A/A \\
\midrule
Equal cost, 32 rows & 5.594 & 1.964 & 2.849 & 1.006 & 0.999 \\
Equal cost, 128 rows & 75.118 & 8.017 & 9.064 & 1.012 & 1.001 \\
Equal cost, 512 rows & 1164.347 & 33.169 & 35.563 & 1.007 & 1.025 \\
Distinct costs, 128 rows & 74.723 & 74.804 & 0.985 & 1.004 & 0.999 \\
24 random 24-row queries & 111.174 & 130.429 & 0.858 & 1.001 & 1.069 \\
Circle, 32 crossings & 1683.749 & 429.510 & 3.914 & 1.000 & 1.015 \\
Circle, 4 crossings, binary & 49.438 & 27.429 & 1.811 & 1.008 & 0.988 \\
\bottomrule
\end{tabular}

\end{center}

On the equal-cost family, complete endpoints improve from $2.85\times$
at 32 rows to $35.56\times$ at 512 rows, consistent with the proved
residual-search improvement. The largest median falls from 1.164 seconds
to 0.033 seconds. The distinct-cost control is close to parity, with a
paired ratio of 0.985 after the bounded handoff. The random control is
about 17 percent slower by its paired ratio; its new-arm A/A ratio of
1.069 also shows visible variation. This is consistent with the adverse
guard-work result above, and prevents a claim of uniformly faster span
optimization.

The two source-derived optimization endpoints improve $3.91\times$ and
$1.81\times$. The 32-crossing case includes source construction and
cohomology, falling from 1.684 to 0.430 seconds. This is a timing of its
optimization query: the recognizer normally succeeds on its raw surface
and does not run that optimization.

The second experiment starts from the input PD and times the complete
recognizer, including its earlier setup, native discovery, exact fallback
when required, an additional independent replay of any returned native
proof and certificate/coordinate hashing. Earlier shortcuts are disabled
to make the native stage reachable. Both arms use no local seed cap and
retain the recognizer's twelve-second global deadline. The late-tree
case requests 24 trials; the others use four. This mode completes 120
measured calls and 24 warm-ups in 11.591 seconds.

\begin{center}
\small\begin{tabular}{lrrrrr}
\toprule
Input & Old (ms) & New (ms) & Old/new & Old A/A & New A/A \\
\midrule
Raw circle & 63.035 & 61.919 & 1.034 & 1.032 & 0.995 \\
Optimized positive & 128.244 & 89.515 & 1.459 & 0.956 & 0.974 \\
Early tree positive & 37.066 & 38.425 & 0.999 & 1.000 & 1.005 \\
Late tree positive & 178.053 & 139.721 & 1.280 & 0.964 & 1.049 \\
Genus-one miss & 57.538 & 41.679 & 1.378 & 0.997 & 0.986 \\
Trefoil & 55.427 & 39.291 & 1.404 & 0.976 & 1.000 \\
\bottomrule
\end{tabular}

\end{center}

The four cases reaching span optimization improve $1.28$--$1.46\times$.
They preserve the optimized and late-tree positive proofs' conclusions,
and preserve the genus-one unknot miss and trefoil's exact fallback
verdicts. The early tree positive has a paired ratio of 0.999 and avoids
flow in both arms. The raw-circle control's 1.034 ratio is comparable to
its old-arm A/A ratio of 1.032; no algorithmic gain is claimed for that
unchanged branch. These are selected native-stage timings, not evidence
of improved ordinary portfolio performance or a general recognition bound.
